\section*{Mòdul de representació intermitja (IR)}
Aquest mòdul és l'encarregat de traduir la sortida del parser, que depèn tant del llenguatge analitzat com de la implementació concreta del parser, a una representació intermèdia (IR) independent d'aquests factors. El seu objectiu és permetre que els mòduls d'anàlisi posteriors accedeixin a la informació del programa mitjançant una estructura comuna, sense haver de realitzar adaptacions específiques per a cada llenguatge o parser. Aquesta idea es reflecteix en el model actual, on els nodes del llenguatge d'origen es converteixen en instàncies de \texttt{ParsedElement} amb un \texttt{ElementKind} normalitzat.

La traducció es basa en el principi de generalització dels components fonamentals dels llenguatges de programació. En lloc d'exposar directament l'\textit{AST} del parser, el sistema defineix categories bàsiques com imports, declaracions, paràmetres, invocacions, expressions i accedors, i les reutilitza per a diferents construccions sintàctiques. Això permet unir conceptes equivalents encara que el llenguatge origen els presenti amb notació diferent.

Per exemple, les declaracions de tipus, els mètodes o els membres d'una classe poden ser representats mitjançant una estructura comuna que conserva la informació rellevant per a l'anàlisi posterior. La conversió s'implementa mitjançant \texttt{ElementTranslator}, que assigna cada tipus de node a una classe del model comú, i la classe \texttt{CSharpElementMapper}, que adapta els nodes de Roslyn a aquests tipus. Tal com es mostra a la figura~\ref{fig:parser_module}, aquesta normalització és el punt de connexió entre el parser i la capa d'anàlisi.

Aquesta generalització té com a objectiu proporcionar una representació independent del llenguatge d'origen, mantenint la informació estructural i semàntica necessària perquè els mòduls d'anàlisi posteriors puguin operar-hi sense dependre del parser utilitzat ni de la seva implementació interna.

